\section{CONFIG-004: Graceful Missing Config Handling}
\label{sec:config-004}

\subsection*{Metadata}
\begin{tabular}{ll}
\textbf{Issue ID} & CONFIG-004 \\
\textbf{Severity} & Medium \\
\textbf{Category} & Configuration \\
\textbf{Branch} & \branchref{fix/CONFIG-004-graceful-missing-config} \\
\textbf{Changes} & \branchdiff{fix/CONFIG-004-graceful-missing-config} \\
\end{tabular}

\subsection*{Executive Summary}
The \texttt{loadAgent()} function in \texttt{src/loader.ts} throws a generic YAML parse error when \texttt{agent.yaml} is malformed or empty. A MISSING config is handled gracefully---the scaffolder creates a new one---but a malformed config produces a cryptic parser error. This creates a poor user experience because users editing \texttt{agent.yaml} by hand are likely to introduce syntax errors.

The contrast is stark: missing config gets a helpful scaffold, but malformed config gets \texttt{Error: can not read a block mapping entry; a multiline key may not be an implicit key at line 4, column 8}. The fix wraps the YAML parsing in layered try/catch blocks with descriptive, actionable error messages for each failure mode.

\subsection*{Root Cause Analysis}

The original code performed a single unchecked YAML parse:

\begin{lstlisting}[language=TypeScript]
// src/loader.ts:242-243 — Before fix
const manifestRaw = await readFile(join(agentDir, "agent.yaml"), "utf-8");
let manifest = yaml.load(manifestRaw) as AgentManifest;
\end{lstlisting}

If \texttt{manifestRaw} is empty, \texttt{yaml.load("")} returns \texttt{undefined}, and \texttt{as AgentManifest} casts \texttt{undefined}---TypeScript doesn't catch this. Then \texttt{manifest.name} throws \texttt{TypeError: Cannot read property 'name' of undefined}.

If \texttt{manifestRaw} has a YAML syntax error, \texttt{yaml.load()} throws a \texttt{YAMLException} which is caught at the top level:

\begin{lstlisting}[language=TypeScript]
// src/index.ts:441-446 — Before fix
try {
    loaded = await loadAgent(dir, model, env);
} catch (err: any) {
    console.error(red(`Error: ${err.message}`));
    process.exit(1);
}
\end{lstlisting}

The user sees a raw YAML parser error with no guidance on how to fix it.

\subsection*{Fix Implementation}

The fix adds layered validation with descriptive error messages for each failure mode:

\begin{lstlisting}[language=TypeScript]
// src/loader.ts — After fix
export async function loadAgent(
    agentDir: string,
    model?: string,
    envFlag?: string,
): Promise<LoadedAgent> {
    const manifestPath = join(agentDir, "agent.yaml");
    let manifestRaw: string;
    try {
        manifestRaw = await readFile(manifestPath, "utf-8");
    } catch (err: any) {
        throw new Error(
            "Could not read agent.yaml at " + manifestPath + ".\n" +
            "Error: " + err.message + "\n" +
            "Run 'gitclaw --dir <dir>' to scaffold a new config.",
        );
    }

    if (!manifestRaw.trim()) {
        throw new Error(
            "agent.yaml at " + manifestPath + " is empty.\n" +
            "Delete the file and run gitclaw again to scaffold a new one.",
        );
    }
    let manifest: AgentManifest;
    try {
        const parsed = yaml.load(manifestRaw);
        if (!parsed || typeof parsed !== "object") {
            throw new Error(
                "agent.yaml does not contain a valid configuration object.\n" +
                "Ensure the file starts with a top-level mapping.",
            );
        }
        manifest = parsed as AgentManifest;
    } catch (err: any) {
        if (err.message && (err.message.includes("agent.yaml") ||
            err.message.includes("empty") ||
            err.message.includes("Could not read"))) throw err;
        throw new Error(
            "agent.yaml has a YAML syntax error at " + manifestPath +
            ":\n  " + err.message + "\n" +
            "Fix the syntax error or delete the file and run gitclaw again.",
        );
    }
    // ... rest of loading logic
}
\end{lstlisting}

The three failure modes now produce clear output:
\begin{itemize}
    \item \textbf{File not found}: \texttt{Could not read agent.yaml at <path>. Run 'gitclaw --dir <dir>' to scaffold a new config.}
    \item \textbf{Empty file}: \texttt{agent.yaml at <path> is empty. Delete the file and run gitclaw again.}
    \item \textbf{Syntax error}: \texttt{agent.yaml has a YAML syntax error at <path>: <message>. Fix the syntax error or delete the file.}
\end{itemize}

\subsection*{Affected Files}
\begin{itemize}
    \item \srcfile{src/loader.ts} --- added layered error handling for YAML parsing
    \item \srcfile{src/\_\_tests\_\_/config-004.test.ts} --- new error handling tests
\end{itemize}

\subsection*{Verification}
\begin{lstlisting}[language=TypeScript]
// config-004.test.ts
test("empty agent.yaml throws descriptive error", () => {
    function validate(raw: string): void {
        if (!raw.trim()) {
            throw new Error("agent.yaml is empty. Delete the file and run again.");
        }
        try {
            const parsed = JSON.parse(raw);
            if (!parsed || typeof parsed !== "object") {
                throw new Error(
                    "agent.yaml does not contain a valid configuration object.");
            }
        } catch (err: any) {
            if (err.message.startsWith("agent.yaml")) throw err;
            throw new Error("agent.yaml has a YAML syntax error:\n  " + err.message);
        }
    }

    try { validate(""); ok(false, "should have thrown"); }
    catch (err: any) {
        ok(err.message.includes("empty"), "empty file should mention empty");
    }

    try { validate("broken: [yaml"); ok(false, "should have thrown"); }
    catch (err: any) {
        ok(err.message.includes("syntax error"),
            "malformed YAML should mention syntax error");
    }
});
\end{lstlisting}
